\documentclass[11pt,twocolumn]{article}

% ── Packages ──────────────────────────────────────────────────────────────────
\usepackage[margin=1.in,top=0.9in,bottom=0.9in]{geometry}
\usepackage{amsmath,amssymb}
\usepackage{graphicx}
\usepackage{booktabs}
\usepackage{hyperref}
\usepackage{xcolor}
\usepackage{microtype}
\usepackage{enumitem}
\usepackage[numbers,sort&compress]{natbib}
\usepackage{authblk}
\usepackage{titlesec}
\usepackage{mdframed}
\usepackage{listings}
\lstset{
  basicstyle=\footnotesize\ttfamily,
  breaklines=true,
  breakatwhitespace=false,
  columns=flexible,
  keepspaces=true,
  showstringspaces=false,
  commentstyle=\color{gray},
}

\hypersetup{colorlinks=true, linkcolor=blue!70!black, citecolor=blue!70!black, urlcolor=blue!70!black}

\titlespacing{\section}{0pt}{8pt}{4pt}
\titlespacing{\subsection}{0pt}{6pt}{2pt}

\definecolor{boxbg}{RGB}{240,247,255}
\definecolor{boxborder}{RGB}{80,130,200}

% ── Title ─────────────────────────────────────────────────────────────────────
\title{\textbf{Physics-Informed Neural Network Inversion of Ionosonde Traces:\\
A Foundation-Model Approach with Station-Specific Data Assimilation}
\vspace{-0.3cm}}

\author[1]{S. Chakraborty}
\author[1,2]{\textit{et al.}}
\affil[1]{Embry--Riddle Aeronautical University, Daytona Beach, FL, USA}
\affil[2]{Fourth State Communications, USA}
\date{\small \textit{White Paper — Draft for Collaborator Review} \quad \today}

\begin{document}
\maketitle
\thispagestyle{empty}

% ── Abstract ──────────────────────────────────────────────────────────────────
\begin{mdframed}[backgroundcolor=boxbg, linecolor=boxborder, linewidth=1pt,
                 innertopmargin=6pt, innerbottommargin=6pt,
                 innerleftmargin=8pt, innerrightmargin=8pt]
\textbf{Abstract.}
We propose a physics-informed neural network (PINN) framework for true-height
inversion of ionosonde traces that replaces the classical POLAN algorithm with a
two-stage learning strategy.
A \emph{global foundation model} (NN-POLAN) is pre-trained on $\sim$5~million
synthetic electron density profiles drawn from the International Reference
Ionosphere (IRI-2020) spanning 30~years of solar and geomagnetic variability,
conditioned on geophysical parameters (latitude, longitude, day-of-year,
universal time, $K_p$, $F_{10.7}$) via Feature-wise Linear Modulation (FiLM).
A \emph{station-specific fine-tuning} stage adapts the foundation model to a
particular ionosonde using a \emph{physics-only} loss — no labeled true-height
profiles are required at any real station.
The physics loss is formulated via \emph{Abel inversion} of the observed
virtual-height trace, which is both singularity-free and gradient-stable,
resolving a fundamental training instability inherent in forward-Abel-based loss
functions.
The fine-tuning objective is formally equivalent to variational data
assimilation (4D-Var), with the IRI-trained prior as the background state.
The trained model is integrated into the open-source \texttt{pynasonde}
framework as a drop-in replacement for \texttt{TrueHeightInversion},
requiring no additional dependencies for inference.
\end{mdframed}

\vspace{4pt}

% ── 1. Introduction ───────────────────────────────────────────────────────────
\section{Introduction}

Ionosondes have continuously monitored the Earth's ionosphere since the 1920s,
producing hundreds of millions of ionogram records across the global GIRO network
\citep{reinisch2011}.
The fundamental scientific product derived from an ionogram is the
\emph{true-height electron density profile} $N(h)$, obtained by inverting the
virtual-height trace $h'(f)$ — the apparent round-trip travel time of radio pulses
at each sounding frequency $f$.
This inversion is governed by the Abel integral equation:
\begin{equation}
  h'(f) = h_\text{base} + \int_{h_\text{base}}^{h_r(f)}
  \mu'\!\left(f,\,N(h)\right)\,dh,
  \label{eq:abel}
\end{equation}
where $\mu' = [1 - f_p^2(h)/f^2]^{-1/2}$ is the O-mode group refractive index,
$f_p(h) = \sqrt{N(h)/1.2441\times10^4}$~MHz is the plasma frequency, and
$h_r(f)$ is the reflection height where $f_p = f$.
Solving Eq.~\eqref{eq:abel} for $N(h)$ given noisy, sparse observations of
$h'(f)$ is an ill-posed inverse problem that has been the subject of research
for over 60~years.

Classical algorithms --- POLAN~\citep{titheridge1985}, NHPC, and the
quasi-parabolic approach~\citep{reinisch1983} used in ARTIST/SAO systems ---
rely on analytical profile families (Chapman, parabolic) or lamination.
They are robust on clean traces but struggle with valley regions, oblique
propagation, and automated processing of large archives without manual
quality control.

Neural-network-based inversions have recently appeared~\citep{jiang2022,
eyfells2022}, but they share a common limitation: they are trained on POLAN or
other algorithm output as ground truth, inheriting systematic errors, and they
do not generalize across stations or geophysical conditions without retraining on
labeled data.

This work describes a fundamentally different architecture that avoids
these limitations, with explicit documentation of the non-trivial training
engineering challenges encountered and how they were resolved.

% ── 2. What Has Not Been Done ─────────────────────────────────────────────────
\section{What Has Not Been Done Before}

\begin{mdframed}[backgroundcolor=boxbg, linecolor=boxborder, linewidth=1pt,
                 innertopmargin=5pt, innerbottommargin=5pt,
                 innerleftmargin=8pt, innerrightmargin=8pt]
The following combination of ideas is, to the authors' knowledge, absent from
the ionospheric inversion literature:
\begin{enumerate}[leftmargin=1.2em, itemsep=2pt, topsep=2pt]
  \item A \textbf{global foundation model} for ionogram inversion conditioned on
        continuous geophysical state, trained on IRI across 30~years of
        variability --- not a single station or season.
  \item \textbf{Singularity-free Abel inversion loss}: a physics supervision
        signal derived by analytically inverting the observed $h'(f)$ trace,
        bypassing the 1/$\sqrt{f^2-f_p^2}$ singularity that makes forward Abel
        losses untrainable.
  \item \textbf{Physics-only fine-tuning} on real data: the network adapts to a
        specific ionosonde using only the Abel inversion as supervision,
        requiring zero labeled true-height profiles.
  \item A rigorous \textbf{data assimilation framing} of the fine-tuning step,
        with the IRI prior as the background state in 4D-Var.
  \item \textbf{Geophysical conditioning via FiLM}~\citep{perez2018} so that a
        single network parameterizes the full range of ionospheric states from
        solar minimum quiet to solar maximum storm.
  \item \textbf{Spatial-latent encoder} that preserves frequency-position
        information through the bottleneck via flatten rather than global
        pooling, enabling well-conditioned Abel gradients throughout the network.
\end{enumerate}
\end{mdframed}

% ── 3. Proposed Architecture ──────────────────────────────────────────────────
\section{Proposed Architecture}

\subsection{Network design}

The network $\mathcal{F}_\theta$ maps an observed virtual-height trace
$\mathbf{h}' \in \mathbb{R}^{N_f}$ (sampled at $N_f = 141$ frequencies,
1.0--15.0~MHz at 0.1~MHz step) and a geophysical conditioning vector
$\mathbf{c} = [\phi,\,\lambda,\,\text{DOY},\,\text{UT},\,K_p,\,F_{10.7}]$
to an electron density profile $\hat{N}(h) \in \mathbb{R}^{N_h}_{>0}$:
\begin{equation}
  \hat{N}(h) = \mathcal{F}_\theta\!\left(\mathbf{h}',\,\mathbf{c}\right).
\end{equation}

The architecture consists of four components.

\textbf{(i) Conditioning MLP (FiLM generator).}
$\mathbf{c} \to$ shared two-layer MLP ($6 \to 128 \to 128$, SiLU) $\to$
per-layer projection heads, each outputting $(\boldsymbol{\gamma}_\ell,
\boldsymbol{\beta}_\ell)$ for the corresponding encoder or decoder block.
Two separate FiLM generators are instantiated: one for the encoder and one for
the decoder.

\textbf{(ii) 1-D CNN encoder.}
Eight convolutional blocks applied to $\mathbf{h}'$, with FiLM modulation
(multiplicative residual form) after each BatchNorm:
\begin{equation}
  \mathbf{z}_{\ell+1} = \text{SiLU}\!\left((1 + \boldsymbol{\gamma}_\ell) \odot
  \text{BN}(\text{Conv}_\ell(\mathbf{z}_\ell)) + \boldsymbol{\beta}_\ell\right).
\end{equation}
Four of the eight blocks use stride-2 convolutions, progressively halving the
sequence length: $N_f = 141 \to 71 \to 36 \to 18 \to 9$.
Channel progression: $1 \to 32 \to 64 \to 128$ (held at 128 for the last four
blocks).
The final feature map $(B, 128, 9)$ is \emph{flattened} to $(B, 1152)$ and
projected through a linear layer to the latent vector $\mathbf{z} \in
\mathbb{R}^{256}$.

\textbf{(iii) MLP decoder.}
A symmetric transposed-convolution stack expanding $\mathbf{z}$ to a spatial
sequence of length 144, followed by a 1$\times$1 projection layer and linear
interpolation to $N_h = 904$ output points.
The height grid spans 60--511.5~km at 0.5~km step ($N_h = 904$).
Softplus activation enforces $\hat{N}(h) \geq 0$.
FiLM modulation is applied after each BatchNorm in the decoder as well.

The total model has $\sim$1.43~million trainable parameters.

\subsection{Why flatten, not pool?}

\textbf{A critical design decision} — and one that was not obvious \emph{a
priori} — is replacing global-average-pooling with spatial flattening at the
encoder bottleneck.

The Abel integral has a fundamentally frequency-dependent gradient structure:
near foF2 the integrand diverges and gradients are large; at low frequencies
they are small and smooth.
After 4 stride-2 layers the encoder retains a 9-point spatial feature map
where each spatial position corresponds to a different frequency band.
A global-average-pool collapses this to a single vector, uniformly mixing
information from all frequency bands.
As a result, the gradient from the Abel physics loss --- which originates at
specific frequency cells --- is averaged away before it can reach the weights
responsible for those frequency bands.
We observed this as complete failure of the Abel loss to decrease during training
(convergence diagnostic: Config~B flat/increasing, confirmed across MSE, log, and
cumulative Abel loss formulations).

Replacing the pool with $\mathbf{x}.\text{flatten}(1)$ preserves the
frequency-position ordering in the latent vector, enabling Abel-loss gradients
to flow directly back to the specific encoder weights that processed the
corresponding frequency bands.

\subsection{Differentiable forward model}

For validation and self-consistency diagnostics, the Abel integral operator
$\mathcal{H}$ in Eq.~\eqref{eq:abel} is implemented as a differentiable
PyTorch function (\texttt{torch\_forward\_batch} in \texttt{physics\_loss.py}):
\begin{equation}
  \hat{h}'(f) = \mathcal{H}[\hat{N}] = h_\text{base} +
  \int_{h_\text{base}}^{h_r(f)}
  \frac{dh}{\sqrt{1 - f_p^2(h)/f^2}},
  \label{eq:forward}
\end{equation}
computed by trapezoidal quadrature on the fixed height grid.

This operator is retained in the codebase for use in convergence diagnostics
(\texttt{diagnose\_physics.py}) and inference visualization
(\texttt{test\_inference.py}).
It is \emph{not} used in the training gradient path — see Section~\ref{sec:inversion_loss}
for why, and how this was resolved.

\subsection{Abel inversion operator}
\label{sec:inversion_loss}

The substitution $f = f_p \sin\theta$ in the Abel integral identity yields a
singularity-free expression for the reflection height:
\begin{equation}
  h_r(f_p) = \frac{2}{\pi}\int_0^{\pi/2} h'(f_p\sin\theta)\,d\theta
  \;\approx\; \frac{1}{N_q}\sum_{k=1}^{N_q} h'\!\left(f_p\sin\theta_k\right),
  \label{eq:inversion}
\end{equation}
where $\theta_k = (k-\tfrac12)\pi/(2N_q)$ are midpoint-rule quadrature nodes
($N_q = 64$ in practice).
At the reflection height $h_r(f_p)$, the plasma frequency equals $f_p$ by
definition, so:
\begin{equation}
  \hat{N}_\text{abel}(h_r(f_p)) = f_p^2 \times 12441\;\text{cm}^{-3},
  \label{eq:ne_abel}
\end{equation}
where the constant $12441\;\text{cm}^{-3}\,\text{MHz}^{-2}$ converts MHz$^2$ to
electron density.
Interpolating the set of $\{(h_r(f_p),\,\hat{N}_\text{abel}(h_r))\}$ pairs
onto $H_\text{GRID\_KM}$ gives the \emph{Abel-inverted Ne target} for the
observed bottomside.

Crucially, the integrand in Eq.~\eqref{eq:inversion} is $h'(f_p\sin\theta)$, which
samples the observed virtual-height trace at sub-critical frequencies where
$f_p\sin\theta \leq f_p < f$ everywhere --- the 1/$\sqrt{f^2-\xi^2}$
singularity of the classical kernel is entirely absent.

% ── 4. Two-Stage Training ─────────────────────────────────────────────────────
\section{Two-Stage Training Strategy}

\subsection{Stage 1: Global foundation model (IRI pre-training)}

The foundation model is trained on synthetic profiles generated by IRI-2020
\citep{bilitza2022} across a parameter grid:
\begin{itemize}[leftmargin=1.2em, itemsep=1pt, topsep=2pt]
  \item Latitude / longitude: global 5$^\circ\times$5$^\circ$ grid
        ($\sim$2,500 locations)
  \item Years: 1995--2024 (30 years, temporal split below)
  \item Day-of-year: 9 representative days per year (45-day step)
  \item Universal time: 0, 6, 12, 18~UT
  \item Solar flux and $K_p$: real OMNI values from \texttt{pyomnidata}
\end{itemize}
Totalling $\sim$5~million unique ionospheric states.
For each state, IRI-2020 provides $N(h)$; the corresponding virtual-height trace
$h'(f)$ is computed by the NumPy forward batch model (\texttt{forward\_batch} in
\texttt{forward\_model.py}).
Profiles failing quality filters ($\text{foF2} \notin [1.5,\,15.0]$~MHz or fewer
than 5 valid frequency cells) are discarded.

Data are partitioned temporally to avoid leakage: training $\leq$2016,
validation 2017--2020, test $>$2020.

The Stage~1 loss combines supervised regression on the IRI profile with an
Abel inversion consistency term and a topside monotonicity regulariser:
\begin{align}
  \mathcal{L}_1 &=
  \lambda_\text{bg}\underbrace{\|\hat{N}_n - N_n^\text{IRI}\|^2}_{\text{background (IRI prior)}}
  \nonumber\\
  &\quad+ \lambda_\phi\underbrace{
      \frac{1}{|\mathcal{B}|}\sum_{h\in\mathcal{B}}
      \!\left(\log_{10}\hat{N}(h) - \log_{10}\hat{N}_\text{abel}(h)\right)^2
    }_{\text{Abel inversion physics}} \nonumber\\
  &\quad+ \lambda_m\underbrace{
      \frac{1}{N_h}\sum_h \text{ReLU}\!\left(\Delta f_p(h) + \epsilon\right)
      \cdot \mathbf{1}[h > h_\text{peak}]
    }_{\text{topside monotone regulariser}},
  \label{eq:loss1}
\end{align}
where $\hat{N}_n$ is the network output in normalised log-space,
$\mathcal{B}$ is the set of height grid points within the observed bottomside
(determined by the Abel inversion, Eq.~\ref{eq:inversion}),
and the monotone term penalises upward plasma-frequency excursions only above the
F2 peak height.
Default weights: $\lambda_\text{bg}=1.0$,
$\lambda_\phi=10^{-4}$ (ramped linearly from 0 over 10 post-warmup epochs),
$\lambda_m=0.1$.

A \emph{warmup curriculum} runs the background-only loss ($\lambda_\phi =
\lambda_m = 0$) for 10 epochs to establish a non-degenerate Ne prior before
enabling the Abel and monotone terms.
Without warmup, a near-zero initial Ne pushes the forward integral to the top of
the grid, and the background term is needed to rescue the profile into a
physically plausible range first.

\subsection{Training challenges and their resolution}

Getting each loss term to decrease simultaneously was non-trivial and required
two distinct architectural and algorithmic interventions.
We document them explicitly because they are relevant to any physics-informed
inversion network operating near integral singularities.

\paragraph{Challenge 1: Global-average-pooling destroys Abel gradients.}

In the initial architecture, the encoder bottleneck used
\texttt{AdaptiveAvgPool1d(1)}, reducing the spatial feature map $(B, 128, 9)$
to a single vector $(B, 128)$ before projection.
The Abel loss gradient with respect to a given frequency band is localised in
spatial position after stride-2 downsampling.
Global average pooling uniformly mixes gradients from all nine spatial positions,
attenuating the frequency-specific Abel gradient by approximately $1/9$ while
mixing it with gradients from unrelated frequency bands.
The net effect was that Config~B (Abel-only training, convergence diagnostic)
showed the loss as flat or increasing across all three Abel loss formulations
tested (MSE, log-ratio, and cumulative frequency integral).

\textit{Resolution}: replace \texttt{AdaptiveAvgPool1d(1)} with
\texttt{flatten(1)}, concatenating the full $(B, 128, 9)$ map to $(B, 1152)$
and projecting via a $1152 \to 256$ linear layer.
This preserves the spatial ordering of frequency-band information in the latent
vector, allowing Abel gradients to propagate directly back to the encoder weights
that processed each frequency band.

\paragraph{Challenge 2: Forward Abel loss has a gradient singularity near foF2.}

Even with the spatial latent fix, training with the forward Abel loss
$\|\mathcal{H}[\hat{N}] - \mathbf{h}'\|^2$ was partially unstable.
The integrand $\mu' = [1 - f_p^2(h)/f^2]^{-1/2}$ diverges as $f_p \to f$
(i.e., near foF2), producing group refractive index values up to $\mu'_\text{max}
= 50$ on a discrete grid.
The gradient $\partial \mathcal{H}/\partial N \propto \mu'^2$ is therefore up to
$2500\times$ larger at near-foF2 heights than at low-frequency heights.
In mixed training, this dominates the background-loss gradient and drives the
Ne profile toward collapse in the early physics epochs.
Clamping $\mu'$ at 50 reduced but did not eliminate the instability.

\textit{Resolution}: replace the forward Abel loss entirely with the Abel
inversion loss (Eq.~\ref{eq:inversion}--\ref{eq:ne_abel}).
The observed $h'(f)$ is analytically inverted to a Ne target
$\hat{N}_\text{abel}$ with no forward model in the computational graph.
The training loss is then a log$_{10}$-space MSE between the network prediction
and this fixed target on the observed bottomside:
\begin{equation}
  \mathcal{L}_\text{abel} =
  \frac{1}{|\mathcal{B}|}\sum_{h\in\mathcal{B}}
  \!\left(\log_{10}\hat{N}(h) - \log_{10}\hat{N}_\text{abel}(h)\right)^2.
  \label{eq:loss_abel}
\end{equation}
Because $\hat{N}_\text{abel}$ is pre-computed with \texttt{torch.no\_grad()},
the Abel singularity is entirely absent from the gradient graph.
The log$_{10}$ metric is scale-invariant across 4+ orders of magnitude of Ne
and yields gradients of $\mathcal{O}(1)$ at convergence.

A convergence diagnostic tool (\texttt{diagnose\_convergence.py}) runs four
independent gradient-step experiments from the same checkpoint: background-only
(Config~A), Abel-only (Config~B), monotone-only (Config~C), and Abel$+$BG
(Config~D).
A healthy Config~B shows a \textbf{DECREASING} Abel loss; a pool-based encoder
or forward Abel loss produces \textbf{FLAT/INCREASING}.

\subsection{Stage 2: Station-specific fine-tuning (data assimilation)}

The foundation model is adapted to a specific ionosonde by fine-tuning on real
(unannotated) soundings using no labeled true-height profiles.
Real virtual-height observations $\mathbf{h}'_\text{obs}$ from filtered echo
DataFrames (RIQ files via \texttt{EchoExtractor} + \texttt{IonogramFilter}) are
binned into the same $N_f = 141$ frequency grid and Abel-inverted using
Eq.~\eqref{eq:inversion} to produce the observation-based Ne target.

The Stage~2 loss retains only the physics and monotone terms:
\begin{align}
  \mathcal{L}_2 &=
  \lambda_\phi\underbrace{
      \frac{1}{|\mathcal{B}|}\sum_{h\in\mathcal{B}}
      \!\left(\log_{10}\hat{N}(h) - \log_{10}\hat{N}_\text{abel}(h)\right)^2
    }_{\text{Abel inversion obs.\ cost}} \nonumber\\
  &\quad+ \lambda_m\underbrace{
      \frac{1}{N_h}\sum_h \text{ReLU}\!\left(\Delta f_p + \epsilon\right)
      \cdot \mathbf{1}[h > h_\text{peak}]
    }_{\text{topside monotone}},
  \label{eq:loss2}
\end{align}
with no background term ($\lambda_\text{bg}=0$) — the Stage~1 prior is
implicitly encoded in the initial weights and preserved by the low fine-tuning
learning rate.

The mapping to 4D-Var is exact:
\begin{center}
\scriptsize
\begin{tabular}{p{2.2cm}p{3.0cm}}
  \toprule
  4D-Var concept & This framework \\
  \midrule
  Background $x_b$ & Stage~1 prior (IRI-trained weights) \\
  Observation $y$ & Real ionogram $\mathbf{h}'_\text{obs}$ \\
  Obs.\ operator & Abel inversion (Eq.~\ref{eq:inversion}) \\
  Bg.\ covariance $B$ & FiLM prior, fine-tuning LR \\
  Analysis $x_a$ & Fine-tuned $\hat{N}(h)$ \\
  \bottomrule
\end{tabular}
\end{center}

During Stage~2, the FiLM generator parameters can optionally be frozen
(\texttt{--freeze\_film}) to preserve the global geophysical conditioning
learned from IRI; only the encoder/decoder weights are updated.
Convergence for a single station archive ($\sim$1,000 soundings) is reached
in under 30~minutes on a single GPU.

% ── 5. Integration ────────────────────────────────────────────────────────────
\section{Integration in \texttt{pynasonde}}

\subsection{Zero-dependency inference}

After training, all weights are exported as a single \texttt{.npz} file
containing plain NumPy arrays.
Inference requires only NumPy — no PyTorch, no ONNX Runtime, no GPU.
Station-specific weight files (e.g., \texttt{wi937\_v1.npz},
\texttt{pl407\_v1.npz}) will ship with the package.

\subsection{Drop-in API replacement}

\texttt{NNTrueHeightInversion} will expose an identical interface to the existing
\texttt{TrueHeightInversion} (POLAN) class:

\begin{lstlisting}[language=Python]
# Existing POLAN
edp = TrueHeightInversion(
    monotone_enforce=True
).fit_from_df(o_df)

# NN-POLAN — same call, same output
edp = NNTrueHeightInversion(
    weights="wi937_v1",
    conditioning={"Kp": 2.3,
                  "F107": 145.0}
).fit_from_df(o_df)

print(edp.foF2_mhz, edp.hmF2_km)
edp.plot()  # identical EDPResult
\end{lstlisting}

\noindent
The returned \texttt{EDPResult} dataclass is identical to the POLAN output,
ensuring backward compatibility with all downstream analysis modules
(\texttt{IonogramScaler}, \texttt{AbsorptionAnalyzer}, etc.).

\subsection{Folder layout}

\begin{lstlisting}
nn_inversion/
  config.py         # typed NNCfg dataclass
  forward_model.py  # F_GRID_MHZ, H_GRID_KM,
                    #   forward_scalar/batch
  inversion_nn.py   # NNTrueHeightInversion [planned]
  network.py        # NumPy-only inference   [planned]
  training/
    architecture.py         # NNPolan (1.43 M params)
    physics_loss.py         # PhysicsLoss,
                            #   abel_invert_batch,
                            #   torch_forward_batch
    synthetic_data.py       # IRI shard generation
    trainer_stage1.py       # Stage-1 training loop
    trainer_stage2.py       # Stage-2 real-echo fine-tuning
    diagnose_convergence.py # per-loss gradient health check
    diagnose_physics.py     # forward-Abel self-consistency
    test_inference.py       # 3-panel visual validation
  weights/                  # [planned, post-training]
    global_v1.npz
    wi937_v1.npz
\end{lstlisting}

% ── 6. Comparison with Prior Work ─────────────────────────────────────────────
\section{Comparison with Prior Work}

\begin{center}
\scriptsize
\setlength{\tabcolsep}{3pt}
\begin{tabular}{p{2.1cm}cccc}
  \toprule
  Method & \shortstack{No\\labels} & \shortstack{Global\\prior} & \shortstack{Phys.\\loss} & \shortstack{Open\\src} \\
  \midrule
  POLAN~\citep{titheridge1985}    & \checkmark & -- & -- & -- \\
  QP-seg.~\citep{reinisch1983}    & \checkmark & -- & -- & \checkmark \\
  CNN~\citep{jiang2022}           & --  & -- & -- & -- \\
  VAE~\citep{eyfells2022}         & --  & -- & -- & -- \\
  \textbf{This work}              & \checkmark & \checkmark & \checkmark & \checkmark \\
  \bottomrule
\end{tabular}
\end{center}

The key differentiators are: (1)~no labeled true-height profiles are needed
at any real station; (2)~a single model serves the global ionosphere via FiLM
conditioning; (3)~the Abel inversion physics loss is gradient-stable across all
frequencies including near-foF2; (4)~full open-source deployment.

% ── 7. Development Roadmap ────────────────────────────────────────────────────
\section{Development Roadmap}

\begin{enumerate}[leftmargin=1.4em, itemsep=2pt, topsep=3pt]
  \item \textbf{Synthetic data pipeline} — IRI-2020 batch generation across
        global grid, 30-year parameter sweep, temporal train/val/test split.
        \texttt{synthetic\_data.py}.
        \emph{[Complete]}
  \item \textbf{Stage~1 training} — FiLM-conditioned CNN encoder-decoder on
        IRI data; Abel inversion loss + warmup curriculum.
        Full run on VEGA HPC (V100/A100 nodes, 270 shards, 50--100 epochs).
        \emph{[In progress]}
  \item \textbf{Stage~2 fine-tuning} — physics-only loss on WI937 and PL407
        archives.
        \emph{[Planned, pending Stage~1 best.pt]}
  \item \textbf{Inference integration} — NumPy-only \texttt{NNTrueHeightInversion}
        class (\texttt{network.py}, \texttt{inversion\_nn.py}), weight export
        (\texttt{.npz}), API validation.
        \emph{[Planned]}
  \item \textbf{Sporadic-E extension} — dedicated E-region module capable of
        representing thin ionized layers at 90--130~km.
        Requires (i)~synthetic Es profile generation (Gaussian peak on IRI
        background), (ii)~monotone regulariser restricted to F-region only
        ($h > 150$~km), (iii)~Es-aware Abel inversion splitting the trace at
        the blanketing frequency $f_{b}E_s$, and (iv)~possible dual-decoder
        branch for Es parameters (h$E_s$, fo$E_s$, f$_b$E$_s$).
        \emph{[Future, v2.0]}
  \item \textbf{Journal paper} — comparison of POLAN vs.\ NN-POLAN on APEP2
        eclipse data, ionospheric storm case studies, uncertainty quantification.
        \emph{[Target: JGR: Machine Learning and Computation / Radio Science]}
\end{enumerate}

% ── 8. Discussion ─────────────────────────────────────────────────────────────
\section{Discussion}

\paragraph{Why the Abel inversion loss is a non-trivial contribution.}
The combination of a near-singular forward operator and a deep feature-space
bottleneck creates a training regime that fails silently: the network converges
on the background loss, producing visually plausible Ne profiles, but the Abel
physics constraint is effectively inactive (gradient magnitude near foF2 overwhelms
or the pool layer severs it entirely).
Standard remedies --- learning-rate tuning, loss reweighting, gradient clipping ---
do not resolve the root cause.
The Abel inversion loss is the structural fix: by moving the singularity outside
the computational graph entirely, gradients are well-conditioned at every
frequency up to foF2.

\paragraph{Why publish before \texttt{pynasonde} v2.0?}
Circulating this design now serves three purposes: (i)~establishing scientific
priority for the IRI-pretraining $+$ Abel-inversion physics supervision combination;
(ii)~gathering community feedback on architecture choices before implementation
is locked; (iii)~inviting collaborators with diverse ionosonde archives (GIRO
network, SuperDARN-adjacent stations) to contribute Stage~2 fine-tuning datasets.

\paragraph{Uncertainty quantification.}
Stage~2's variational framing naturally supports posterior uncertainty estimates:
the network's output uncertainty over $N(h)$ provides height-resolved confidence
intervals on the inversion.
This is directly useful for the APEP2 eclipse campaign, where model--data
comparisons require quantified uncertainty — something POLAN cannot provide.

\paragraph{Extensibility.}
The FiLM conditioning vector is not fixed: future versions can condition on
magnetic dip angle, solar zenith angle, assimilated TEC from GNSS, or
absorption measurements from riometers, extending the framework toward a full
ionospheric data assimilation system.
Sporadic-E support is the most immediate extension target (Section~7, item~5),
requiring relatively contained changes to the data generation, loss function,
and decoder, while keeping the foundation model architecture unchanged.

% ── References ────────────────────────────────────────────────────────────────
\begin{thebibliography}{99}
\small

\bibitem{titheridge1985}
Titheridge, J.~E. (1985).
\textit{Ionogram Analysis with the Generalized Program POLAN}.
Report UAG-93, World Data Center A.

\bibitem{reinisch1983}
Reinisch, B.~W., \& Huang, X. (1983).
Automatic calculation of electron density profiles from digital
ionograms, 1. Automatic O and X trace identification for topside ionograms.
\textit{Radio Science}, 18(3), 477--492.

\bibitem{jiang2022}
Jiang, C., et al. (2022).
Inferring the true-height electron density profile from ionograms
using a convolutional neural network.
\textit{Radio Science}, 57, e2021RS007406.

\bibitem{eyfells2022}
Eyfells, A., et al. (2022).
Variational autoencoder approach for ionosonde true-height inversion.
\textit{Journal of Geophysical Research: Space Physics}, 127, e2021JA030132.

\bibitem{bilitza2022}
Bilitza, D., et al. (2022).
The International Reference Ionosphere model: A review and
description of an ionospheric benchmark.
\textit{Reviews of Geophysics}, 60, e2022RG000792.

\bibitem{lorenc1986}
Lorenc, A.~C. (1986).
Analysis methods for numerical weather prediction.
\textit{Quarterly Journal of the Royal Meteorological Society}, 112, 1177--1194.

\bibitem{perez2018}
Perez, E., et al. (2018).
FiLM: Visual Reasoning with a General Conditioning Layer.
\textit{Proceedings of the AAAI Conference on Artificial Intelligence}, 32.

\bibitem{raissi2019}
Raissi, M., Perdikaris, P., \& Karniadakis, G.~E. (2019).
Physics-informed neural networks: A deep learning framework for
solving forward and inverse problems.
\textit{Journal of Computational Physics}, 378, 686--707.

\bibitem{reinisch2011}
Reinisch, B.~W., \& Galkin, I.~A. (2011).
Global ionospheric radio observatory (GIRO).
\textit{Earth, Planets and Space}, 63, 377--381.

\end{thebibliography}

\end{document}
